\chapter{نتیجه‌گیری}
\label{ch5}

در این پژوهش، ابتدا مسئله‌ی مسیریابی و یافتن مسیر بهینه در محیط‌های مختلف مورد بررسی قرار گرفت و سپس الگوریتم‌های پایه و پرکاربرد در این حوزه، به‌ویژه الگوریتم \lr{A*}، تحلیل شدند. با توجه به محدودیت‌های موجود در روش‌های کلاسیک، نسخه‌های بهبود‌یافته‌ای از این الگوریتم‌ها معرفی و ارزیابی شدند تا بتوان دقت، سرعت و کیفیت مسیرهای تولیدشده را افزایش داد.

در ادامه، با تمرکز بر ساختار و عملکرد الگوریتم‌های منتخب، رویکردهای مختلفی برای کاهش زمان اجرا و بهبود فرآیند جستجو پیشنهاد شد. نتایج شبیه‌سازی‌ها نشان داد که به‌کارگیری این روش‌ها می‌تواند در بسیاری از سناریوها موجب کاهش تعداد گره‌های بررسی‌شده و در نتیجه افزایش کارایی الگوریتم شود. همچنین مقایسه‌ی نتایج در اندازه‌های مختلف محیط نشان داد که عملکرد الگوریتم پیشنهادی در محیط‌های بزرگ‌تر نیز قابل قبول و مؤثر است.

از سوی دیگر، نتایج به‌دست‌آمده بیانگر آن است که استفاده از ساختارهای کمکی و تکنیک‌های بهبود جستجو می‌تواند تأثیر قابل توجهی در کاهش پیچیدگی عملیاتی داشته باشد. با این حال، همچنان چالش‌هایی مانند وابستگی به نوع محیط، موانع موجود و نحوه‌ی تنظیم پارامترها باقی می‌ماند که می‌تواند در تحقیقات آینده مورد توجه قرار گیرد.

در مجموع، نتایج این پژوهش نشان می‌دهد که الگوریتم‌های بهبود‌یافته‌ی مبتنی بر \lr{A*} می‌توانند گزینه‌ای مناسب برای حل مسائل مسیریابی در کاربردهای مختلف باشند. انتظار می‌رود در پژوهش‌های آتی، با توسعه‌ی بیشتر این روش‌ها و استفاده از داده‌های واقعی‌تر، کارایی و دقت آن‌ها بیش از پیش افزایش یابد.

